\section{Discussion and Conclusion}

In this paper, we have presented a generic approximation and generalization theory and applied it to Gaussian Random Design. Furthermore, our theory is applicable to scenarios where data are sampled from a mixture of distributions, denoted as $P_{\rm data} = \sum_{i} w_i P_i$. Each distribution $P_i$ is assigned a weight $w_i$ and has a low-dimensional support. In such cases, we focus on a subset of distributions with significant weights, while neglecting distributions with small weights. Then the effective Minkowski dimension depends on the data support of the selected distributions. As the sample size grows, including more distributions becomes necessary to achieve higher estimation accuracy. Another example where our theory can be applied is the case of an approximate manifold, where the data are concentrated on a low-dimensional manifold. In this scenario, the effective Minkowski dimension corresponds to the intrinsic dimension of the manifold.

To illustrate the concept of effective dimension in real-world examples, we refer to a study by \citet{pope2021intrinsic}, which investigates the intrinsic dimension of several popular benchmark datasets for deep learning. The intrinsic dimension presented in \citet{pope2021intrinsic} can be seen as an approximate estimate of the Minkowski dimension, as demonstrated in \citet{levina2004maximum,grassberger1983measuring}. In our work, we adopt the methodology employed by \citet{pope2021intrinsic} and utilize generative adversarial networks trained on the ImageNet dataset to generate samples containing varying numbers of daisy images. To estimate the intrinsic dimension of these generated samples, we employ the Maximum Likelihood Estimation (MLE) method, which is achieved by computing the Euclidean distances between each data point and its $k$ nearest neighbors. The obtained results are presented in Figure \ref{figure}, which clearly demonstrates that the intrinsic dimension estimated from a finite sample of images increases as the sample size grows. This finding aligns with our theory that the effective Minkowski dimension is an increasing function of the sample size.
\begin{figure}
    \centering
    \includegraphics[width=0.61\textwidth]{Figures/daisy_25.pdf}
    %\includegraphics[width=0.425\textwidth]{Figures/soap_15.pdf}
    % \includegraphics[width=0.95\linewidth]{images/wordclouds.png}
    % \includegraphics{images/figure1.png}
    \caption{Dimensionality estimates of images obtained using the MLE method with $k$ nearest neighbors under different sample size.}
    \label{figure}
\end{figure}

In conclusion, this paper studies nonparametric regression of functions supported in $\RR^d$ under data distribution with effective Minkowski dimension $p<d$, using deep neural networks. Our results show that the $L^2$ error for the estimation of $f^* \in \cH(\beta,\cD_{\rm data})$ converges in the order of $n^{-2\beta/(2\beta+p)}$. To obtain an $\epsilon$-error for the estimation of $f^*$, the sample complexity scales in the order of $\epsilon^{-(2\beta+p)/\beta}$, which demonstrates that deep neural networks can capture the effective Minkowski dimension $p$ of data distribution. Such results can be viewed as theoretical justifications for the empirical success of deep learning in various real-world applications where data are approximately concentrated on a low-dimensional set.